% Source: source_md/intelligence_part2_combined_ja.md
\section{推論演算子 $F$ と推論層 $L_F$ — 相互拘束・閾値・開放窓}


本章は三つの部分からなる。前半(\S{}2.1〜\S{}2.4)では推論演算子 $F$ を定義し、その作用域として推論層 $L_F$ を確立する。後半(\S{}2.5〜\S{}2.7)では、推論が作用する環境を相互拘束 $H_{ij}$ のパラメータとして捉え、応答的/沈黙的の区別が本体論的範疇ではなく閾値現象であることを示し、閾値移動能力 $\mu_\theta$ を導入する。末尾(\S{}2.8)では、$\mu_\theta$ が三層動力学における開放窓であることを予告する。

\bigskip


\subsection{推論演算子 $F$ の定義}


推論の最小構造は、状態から状態への遷移である:

\[
x_{n+1} = F(x_n)
\]

ここで $x_n$ は仮説・解釈・理解のいずれかの状態を表し、$F$ はそれをより正当化された次状態へ写す演算子である。

この $F$ は VED 理論ツリーにおいて孤立した記号ではない。Bio-IFGT は生命系の状態発展を

\[
\frac{dx}{dt} = F(x;\theta)
\]

と書き、$F$ を「DNA 由来の拘束パラメータ $\theta$ の下での状態発展演算子」とした。本論文の推論演算子 $F$ は、この状態発展演算子の \textbf{知性層における特殊形} である。すなわち、生命系の一般的な状態発展が、第二アトラクター(知性)近傍において推論として現れる。

両者が同一の記号を持つのは偶然ではなく、理論的連続性の現れである。生命の状態発展と知性の推論は、別個の操作ではなく、同一の動力学が異なる層で具現化したものである。記号 $F$ の層をまたいだ再会は、構造が連続していることの証拠として読まれるべきである。

\bigskip


\subsection{推論演算子の四性質}


推論演算子 $F$ は、以下の四つの最小性質によって特徴づけられる。これらは $F$ の具体的実装(演繹的・確率的・統計的)を問わず成り立つ。

\textbf{性質 (F1) — 概念地景内部性}:
$F$ の作用は、既存の概念地景(C-landscape)の内部に留まる。$F$ は新しい次元や新しい概念を無から生成しない。$F$ は既にある地形の上を動く。

\textbf{性質 (F2) — 確率構造}:
$F$ の遷移は決定論的とは限らず、確率構造に従いうる。同一の $x_n$ から複数の $x_{n+1}$ が生じうる。これが後の分岐(\S{}3.7)の基礎となる。

\textbf{性質 (F3) — 連続的反復性}:
$F$ は反復適用される。$x_0 \to x_1 \to x_2 \to \cdots$ という連鎖が推論軌道を形成する。単一の適用ではなく、反復こそが推論の本体である。

\textbf{性質 (F4) — 局所演繹性}:
各遷移 $x_n \to x_{n+1}$ は局所的な正当化関係に従う。大域的な真理保証ではなく、局所的な「次への移行の妥当性」が $F$ を駆動する。

これら四性質は、推論を「正当化を局所的に拡張しながら概念地景上を反復的に動く確率的操作」として規定する。この規定は \S{}3 の非閉包定理の前提(P1)〜(P3)と整合する。

\bigskip


\subsection{推論層 $L_F$ の定義}


推論演算子 $F$ の作用域全体を \textbf{推論層} $L_F$ と呼ぶ。

\textbf{定義 2.1 (推論層 $L_F$)}

\begin{quote}
推論層 $L_F$ は、推論演算子 $F$ が反復適用される状態空間と、その上に $F$ が生成する推論軌道の全体である。
\end{quote}


$L_F$ の最も重要な構造的性質は、\textbf{内部停止条件を持たないこと}である。これは \S{}3 で定理として証明される。直観的に述べれば、$L_F$ は流れの層であり、流れである限り内部で静止しない。

$L_F$ は三層動力学の最下層である。後に導入されるログ層 $L_R$(\S{}5)と創発停止層 $L_E$(\S{}5)は、$L_F$ の運動を基盤として成立する。$L_F$ なしに $L_R$ も $L_E$ もない。

\bigskip


\subsection{Part I との整合 — Physarum・Transformer}


Part I \S{}5 では Physarum(粘菌)と Transformer が知性の領域横断インスタンスとして示された。本論文の枠組みでは、両者はいずれも推論層 $L_F$ を実装している。

\begin{itemize}
\item \textbf{Physarum}: 化学勾配に沿った原形質流動。$\nabla C$ に沿った探索は $F$ の作用そのものである。
\item \textbf{Transformer}: トークン逐次生成。次トークン予測の連鎖は $F$ の反復適用である。
\end{itemize}


両者とも $L_F$ を持つ。しかし両者は $L_R$ と $L_E$ の構造において大きく異なり、さらに $L_F$ のみでは説明されない停止挙動を示す。この差異は \S{}5(三層プロファイル比較)で詳しく分析される。本章ではまず、両者が $L_F$ を共有することを確認するに留める。

ここで重要なのは、$L_F$ の共有が「同じ知性である」ことを意味しないという点である。$L_F$ は知性の必要条件のひとつにすぎず、知性の十分条件ではない。十分条件は三層動力学の全体にある。

\bigskip


\subsection{相互拘束 $H_{ij}$ と推論環境}


推論演算子 $F$ は真空中で作用するのではない。$F$ は環境との相互作用の中で走る。この環境との関係を、VED 基本式の相互拘束 $H_{ij}$ によって捉える。

VED の中心式のひとつは:

\[
C_{ij} = 1 - \exp(-\lambda \rho_i H_{ij})
\]

である。ここで $H_{ij}$ はノード $i$ と $j$ の間の相互拘束、$\rho_i$ は密度的因子、$\lambda$ は結合定数。閉包度 $C_{ij}$ は相互拘束 $H_{ij}$ の関数として決まる。

知性層において、ノード $i$ を推論主体、ノード $j$ を環境(他者・自然・過去・課題)とすれば、両者の閉包度 $C_{ij}$ は相互拘束 $H_{ij}$ のパラメータ位置によって決まる。すなわち、推論主体と環境がどれだけ相互に拘束し合っているかが、その推論がどれだけ閉じうるかを決める。

\bigskip


\subsection{応答的/沈黙的は閾値現象である}


ここから本論文の重要な構造的主張のひとつが導かれる。

推論環境は、しばしば「応答的世界」と「沈黙的世界」に分けられる(\S{}4 で詳述)。応答的世界では推論出力に環境が応答を返し、沈黙的世界では応答が返らない。しかし、この二分は \textbf{本体論的範疇ではない}。

\textbf{主張 2.2 (二分の閾値性)}

\begin{quote}
応答的/沈黙的の区別は、世界が二種類存在することを意味しない。同一の環境が、相互拘束 $H_{ij}$ のパラメータ位置と、推論主体の $(v_{\text{inf}}, r_{\text{res}}, \eta)$ の組によって、応答的にも沈黙的にも現れる。両者は連続的な閾値の両側である。
\end{quote}


具体的に述べれば、同じ自然現象が、ある相互拘束の設定では「応答する世界」(問いかけに兆候が返る)として、別の設定では「沈黙する世界」(何も返らない)として現れる。決定的なのは世界の側の本性ではなく、推論主体と環境の相互拘束パラメータの位置である。

この主張は VED の基底原理の直接の帰結である。VED は静的な範疇ではなく運動を基底に置く。応答性もまた、固定された世界の属性ではなく、相互拘束パラメータ上の位置として動的に決まる。「ここから応答的、ここまで沈黙的」という固定的境界を引くことは、VED の流儀に反する。境界は連続的に移動する閾値である。

\bigskip


\subsection{閾値移動能力 $\mu_\theta$}


応答的/沈黙的が閾値現象であるならば、その閾値を動かす能力が問題になる。

\textbf{定義 2.3 (閾値移動能力 $\mu_\theta$)}

\begin{quote}
閾値移動能力 $\mu_\theta$ は、推論主体が状況に応じて応答性閾値を動かす動力学的能力である。$\mu_\theta$ が高いほど、主体は同一環境を応答的にも沈黙的にも扱い分けることができる。
\end{quote}


ここで $\mu_\theta$ の添字 $\theta$ は、\S{}1.3 で導入した運動可能領域の拘束パラメータ $\theta$ と同じものである。$\mu_\theta$ は単に「閾値を動かす」のではなく、より正確には \textbf{運動可能領域を画定する拘束パラメータ $\theta$ 自体を動かす能力} である。

この点を精密にするため、二種類の運動を区別する。

\textbf{(A) 領域内移動}:
固定された運動可能領域 $\Omega_\theta$ の内部で、推論状態 $x$ が移動する。これは推論演算子 $F$ の通常の作用であり、$\theta$ は動かない。$L_F$ の運動はすべてこの型である。

\textbf{(B) 領域の組み替え}:
運動可能領域 $\Omega_\theta$ の輪郭そのものが変わる。何が到達・維持可能かの地形が組み替わる。これは $\theta$ 自体が動くことを意味する。

$\mu_\theta$ は \textbf{(B) を内発的に起こす能力} である。(A) は $F$ さえあれば起こるが、(B) は $\theta$ を動かす別の能力を必要とする。粘菌や Transformer は (A) を行うが、(B) を内発的には行えない(\S{}5 で詳述)。サピエンス的知性の動力学的特徴は、(B) を内発的に、しかも短い時間スケールで行えることにある。

(B) の有効機能は、単に領域の輪郭を変えることそのものにあるのではない。\S{}5 で示すように、(B) による組み替えは、相互拘束 $H_{ij}$ に沿って \textbf{停止演算子 $E$ を設置できる領域を状態空間上に投影する} 操作である。$L_F$ 単独の状態空間には存在しなかった「ここで止まってよい」という領域が、(B) の投影によって出現する。この投影は意識編で IFGT の第四プリミティブとして導入された projection の、知性層における作用である(\S{}5.5 で詳述)。すなわち $\mu_\theta$ は閾値を動かす能力であると同時に、停止の宿る場所を作る能力でもある。

\subsubsection{$\theta$ 可変性の時間スケール}


生命層と知性層は、$\theta$ を「運動可能領域を画定する拘束パラメータ」として共有する。両層の違いは、$\theta$ の可変性の時間スケールにある。

\begin{longtable}{@{}p{0.18\linewidth}p{0.18\linewidth}p{0.18\linewidth}p{0.18\linewidth}p{0.18\linewidth}@{}}
\toprule
層 & $\theta$ の内容 & 可変性 & 時間スケール & 担い手 \\
\midrule
\endfirsthead
\toprule
層 & $\theta$ の内容 & 可変性 & 時間スケール & 担い手 \\
\midrule
\endhead
生命(Bio-IFGT) & DNA・反応傾向・閾値 & 可変だが超長期 & 世代を跨ぐ & 進化・選択・遺伝 \\
知性(本論文) & 推論操作の拘束・履歴的閾値 & 可変で短期 & 個体内・対話内 & 閾値移動 $\mu_\theta$ \\
\bottomrule
\end{longtable}


生命層では $\theta$(DNA)は個体の生涯ではほぼ固定であり、その可変性は世代を跨ぐ超長期スケールでのみ現れる(進化)。知性層では $\theta$ は個体内で、極端には一回の対話の中で動きうる。同一の $\theta$-可変性原理が、超長期(進化)と短期(認知)という二つの時定数で現れている。

これにより DNA と知性は、別個の現象ではなく、\textbf{同一原理の時間スケール変奏} として統一的に捉えられる。これは VED が時間・空間・生命・意識を「別物」ではなく「同一原理の異なる領域」として再配置してきた操作の、知性層における延長である。

\subsubsection{持続可能な因果推論運用との関係}


閾値移動能力 $\mu_\theta$ は、因果推論を\textbf{能動的・自律的に運用する}ことの条件である。ここで区別を要する。$\mu_\theta$ を持たない知性(粘菌・人工推論システム)も停止しうる——それらは外部条件や環境相互干渉に受動的に従って $L_E$ を成立させる(\S{}5.4)。無限後退から逃れること自体は、$\mu_\theta$ がなくとも、外部依存停止源があれば可能である。

$\mu_\theta$ が可能にするのは、これとは別のことである。$\mu_\theta$ を持つ主体は、環境を固定的に応答的または沈黙的として扱うのではなく、状況に応じて能動的に閾値を動かし、どの領域で因果推論を駆動しどの領域で停止を供給するかを自律的に選べる。これにより、外部条件の変化を待たずに、因果推論を持続的かつ能動的に運用できる。$\mu_\theta$ を持たない主体は停止できるが、その停止条件の組み替えを外部に依存する。$\mu_\theta$ を持つ主体は、その組み替えを個体内で自律的に行う。これがサピエンス的知性における持続可能な因果推論運用の核である。

\bigskip


\subsection{開放窓としての $\mu_\theta$ — 第5節への橋渡し}


ここまでで、推論層 $L_F$(\S{}2.1〜\S{}2.4)と閾値移動能力 $\mu_\theta$(\S{}2.5〜\S{}2.7)を導入した。両者の関係を予告して本章を閉じる。

$L_F$ は内部停止条件を持たない(\S{}3 で証明)。にもかかわらず実際の知性は停止する。この停止はどこから来るのか。

答えの鍵が $\mu_\theta$ である。$\mu_\theta$ は単に閾値を動かす能力にとどまらず、\textbf{$L_F$ から、その運動の痕跡であるログ層 $L_R$ を経由して、創発停止層 $L_E$ へと至る開放窓} である。閾値移動とは、外部相互干渉を推論過程に取り込む動力学的開口部であり、これによって個体内に外部依存型の停止層を作る回路が開く。

すなわち:

\[
L_F \xrightarrow{\;\text{運動の痕跡}\;} L_R \xrightarrow{\;\mu_\theta \text{による開放}\;} L_E
\]

この機構の本格的な展開は \S{}5 で行われる。そこで、ログ層 $L_R$ の定義、創発停止層 $L_E$ の創発、外部停止演算子 $E$ の構造、そして外部参照場 $\Phi_{\text{ext}}$ が個体間 $L_E$ ネットワーク $\Phi_{\text{net}}$ と同型であることが示される。

その前に \S{}3 で、$L_F$ が内部停止できないという構造的事実——本論文の第一中心——を定理として確立する。

\bigskip


\subsection{要約}


本章で確立したのは次の点である:

\begin{enumerate}
\item 推論演算子 $F$ は Bio-IFGT の状態発展演算子 $F(x;\theta)$ の知性層特殊形であり、記号の連続性は理論的連続性を反映する
\item $F$ は四性質(概念地景内部性・確率構造・連続的反復性・局所演繹性)で特徴づけられる
\item $F$ の作用域として推論層 $L_F$ を定義した。$L_F$ は内部停止条件を持たない(\S{}3 で証明)
\item Physarum・Transformer はいずれも $L_F$ を実装するが、$L_F$ の共有は同一の知性を意味しない
\item 応答的/沈黙的の区別は本体論的範疇ではなく、相互拘束 $H_{ij}$ のパラメータ上の閾値現象である
\item 閾値移動能力 $\mu_\theta$ は運動可能領域の拘束パラメータ $\theta$ を内発的に動かす能力(領域の組み替え (B))であり、$F$ の領域内移動 (A) とは区別される
\item $\theta$ 可変性の時間スケールが生命層(超長期)と知性層(短期)を統一する
\item $\mu_\theta$ は $L_F$ から $L_R$ を経由して $L_E$ へ至る開放窓である(\S{}5 で展開)
\end{enumerate}


それでは \S{}3 において、推論層 $L_F$ の内部不動点不在を定理として証明する。

\bigskip
